% tex.web 494's pass_text counts nested conditionals while it skips, and it
% matches cur_cmd -- the MEANING -- so a \let alias of \iffalse is an if_test
% exactly where a spelt-out one is. \newif builds every LaTeX switch as such an
% alias, so a skip that matched the SPELLING counted no nested conditional:
% size10.clo's `\if@compatibility \if@twocolumn ... \else ... \fi \else' took
% the INNER \else for the outer one's, and the outer \else was then `Extra'.
\catcode`\{=1 \catcode`\}=2
\let\ifOUTER\iffalse
\let\ifINNER\iffalse
\ifOUTER
  \ifINNER \message{[A]}\else \message{[B]}\fi
\else
  \message{[C]}
\fi
% The same nesting with the outer test TRUE, so the inner one really runs.
\let\ifOUTER\iftrue
\ifOUTER
  \ifINNER \message{[D]}\else \message{[E]}\fi
\else
  \message{[F]}
\fi
% An aliased \fi has to close the conditional it was reached in, too.
\let\myfi\fi
\let\ifOUTER\iffalse
\ifOUTER \message{[G]}\else \message{[H]}\myfi
\message{[END]}
\end
